\documentclass[11pt,a4paper]{article}
\usepackage[a4paper,margin=2.2cm,top=2cm,bottom=2cm]{geometry}
\usepackage{fontspec}
\setmainfont{Inter}[UprightFont={* Light},BoldFont={* SemiBold},ItalicFont={* Light Italic},BoldItalicFont={* SemiBold Italic}]
\newfontfamily\disp{Inter Display}[UprightFont={* SemiBold}]
\usepackage{amsmath,amssymb,bm}
\usepackage{xcolor,tikz,enumitem,titlesec,graphicx,booktabs,fancyvrb,float,caption}
\usepackage[hidelinks]{hyperref}
\usetikzlibrary{arrows.meta,positioning}
\definecolor{ink}{HTML}{111827}
\definecolor{mute}{HTML}{6B7280}
\definecolor{line}{HTML}{E5E7EB}
\definecolor{acc}{HTML}{0369A1}
\definecolor{soft}{HTML}{F0F9FF}
\color{ink}
\titleformat{\section}{\disp\large\color{ink}}{}{0pt}{}
\titlespacing*{\section}{0pt}{14pt}{5pt}
\titleformat{\subsection}{\small\bfseries\color{acc}}{}{0pt}{}
\titlespacing*{\subsection}{0pt}{9pt}{3pt}
\setlength{\parindent}{0pt}
\setlength{\parskip}{5pt}
\linespread{1.12}
\setlist{nosep,leftmargin=1.2em,itemsep=3pt,label={\color{acc}\textbullet}}
\pagestyle{plain}
\emergencystretch=3em
\sloppy
\captionsetup{font=small,labelfont={bf,color=acc}}

% A visible marker for anything that must come from real hardware runs.
\newcommand{\todo}[1]{\textcolor{acc}{\textbf{[TODO: #1]}}}
\newcommand{\figorti}[3]{\IfFileExists{#1}{\includegraphics[width=#2\linewidth]{#1}}{\colorbox{soft}{\parbox{\dimexpr#2\linewidth-2\fboxsep}{\todo{#3}}}}}
\newcommand{\callout}[1]{\colorbox{soft}{\parbox{\dimexpr\linewidth-2\fboxsep}{#1}}}

\begin{document}

{\small\color{acc}\textbf{SCE 672 \,|\, TERM PROJECT ASSESSMENT 1}}\par\vspace{6pt}
{\disp\fontsize{24}{28}\selectfont RoboMaster EP: Provisioning, Teleop, ROS 2}\par\vspace{2pt}
{\large\color{mute} Keyboard teleoperation, telemetry, and middleware integration}\par\vspace{10pt}
{\small Ameen Ur Rehman (202527430) \quad\textcolor{line}{|}\quad Prof. Aamer Iqbal Bhatti \quad\textcolor{line}{|}\quad KFUPM, 2026}\par
\vspace{6pt}{\color{line}\rule{\linewidth}{0.8pt}}

\section{Problem 1: Host environment and middleware}

\subsection{(i)--(iv) Provisioning snapshots}
Captured from the lab host with the commands shown (raw file: \texttt{logs/system\_diagnostics.txt}).
Kernel, Ubuntu 24.04 LTS, Python 3.12, ROS 2 Jazzy, Gazebo Harmonic (Gazebo Sim 8.x) and NumPy 1.26.4 inside \texttt{\textasciitilde/rm\_venv} are all in place;
\texttt{robomaster} SDK (community fork) and \texttt{robomaster\_ros} were built with \texttt{colcon build --symlink-install}.

{\footnotesize
\VerbatimInput[fontsize=\scriptsize]{../logs/system_diagnostics.txt}}

\subsection{(v) ABI incompatibility analysis}
\textbf{Why \texttt{pip install -U numpy} breaks ROS 2 nodes.}
NumPy 2.0 changed the C-API and ABI. Struct layouts exposed to extensions (for example \texttt{PyArray\_Descr}, whose \texttt{elsize} and \texttt{alignment} fields were widened and moved, and removed legacy members) and the function table behind \texttt{PyArray\_API} were reorganised.
A compiled extension does not call NumPy by name. It indexes into that table and reads those structs at \emph{fixed offsets baked in at build time}.
The Ubuntu/ROS binaries (\texttt{cv\_bridge}'s Boost.Python module, \texttt{rclpy} message converters, OpenCV's \texttt{cv2}) are compiled against NumPy 1.26. If NumPy 2.x is what gets imported at run time, those offsets point at the wrong slots and fields.
The outcome is one of two failure modes: NumPy's guard raises \texttt{ImportError: A module that was compiled using NumPy 1.x cannot be run in NumPy 2.x}, or, where the mismatch goes unchecked, the extension dereferences a wrong pointer and the interpreter dies with a segmentation fault.
Python-level code is unaffected, which is why the failure shows up only at the C boundary (image conversion, point-cloud and array message handling).

\textbf{Why a plain \texttt{-U} triggers it.}
With a venv created using \texttt{--system-site-packages}, an unconstrained upgrade installs NumPy 2.x into the venv's own \texttt{site-packages}. That directory comes first on \texttt{sys.path}, so it \emph{shadows} the apt-provided 1.26.4 that the ROS binaries were built for. The same shadowing happens if a user-level copy sits in \texttt{\textasciitilde/.local}.

\textbf{How the setup isolates and resolves it.}
\begin{itemize}
  \item \texttt{python -m venv --system-site-packages}: the interpreter can see ROS's compiled bindings (\texttt{rclpy}, \texttt{cv\_bridge}) while pip installs stay out of the system Python, which PEP 668 protects.
  \item \texttt{"numpy<2"}: the explicit pin forbids pip from ever resolving to 2.x (the RoboMaster SDK and \texttt{numpy-quaternion} would otherwise pull the latest).
  \item \texttt{rm\_setup.sh}: sources ROS and the workspace overlay, activates the venv, prepends the venv \texttt{site-packages} to \texttt{PYTHONPATH}, and exports \texttt{PYTHONNOUSERSITE=1} so \texttt{\textasciitilde/.local} packages cannot shadow anything. Every shell therefore resolves the same, ABI-consistent NumPy.
\end{itemize}
\textbf{Observation on this host.} \texttt{import numpy} resolves to \texttt{/usr/lib/python3/dist-packages} (1.26.4), because the pin was already satisfied by the system copy and pip therefore installed no second one. The effect is the same: one NumPy, the one ROS was compiled against.

\section{Problem 2: SDK teleoperation and telemetry}

\subsection{(i) Implementation: \texttt{ep\_teleop\_sdk.py}}
A single-file script, importing only \texttt{robomaster} (no ROS 2 imports).
\begin{itemize}
  \item \textbf{Keyboard:} \texttt{curses}, non-blocking. W/S $\to v_x$, A/D $\to v_y$, Q/E $\to \omega_z$. Terminals report presses and auto-repeat but never key release, so each key carries a hold window (0.55\,s for the first press to cover the OS typematic delay, 0.15\,s once repeating). When repeats stop the key counts as released and its contribution drops out. Keys can combine (e.g.\ forward plus strafe).
  \item \textbf{Dispatch:} commands go to \texttt{ep\_robot.chassis.drive\_speed()} at 20\,Hz.
  \item \textbf{Telemetry:} \texttt{sub\_battery\_info}, \texttt{sub\_attitude} and \texttt{sub\_esc} are subscribed at 10\,Hz. A logger writes one CSV row per 100\,ms: battery state of charge, yaw/pitch/roll, four wheel speeds (rpm), the arrival time of the newest ESC packet, and the commanded twist.
  \item \textbf{Failsafe:} Space zeroes $v_x=v_y=\omega_z=0$ immediately (sent in the same loop iteration, not on the next 20\,Hz tick). Releasing keys stops the chassis through the hold window. The \texttt{finally} block sends stop commands and unsubscribes on any exit, including exceptions.
\end{itemize}
SDK sign convention: $x$ forward, $y$ to the \emph{right}, $\omega_z$ \emph{clockwise}. This is the opposite of ROS REP-103 ($y$ left, $\omega_z$ counter-clockwise); the driver handles the conversion for the ROS runs.

\textbf{Logic check without hardware.} The control path was exercised against a stub SDK in a pseudo-terminal: forward hold produced $v_x=0.3$\,m/s, forward plus strafe combined, both dropped to zero 0.45\,s after the last keypress, yaw produced $45^\circ$/s, and Space zeroed all commands within 6\,ms. The CSV logger wrote 10 rows/s. This is a software test only and is \emph{not} the hardware result.

\subsection{(ii) Wired RNDIS run (30\,s)}
Controller switch on Direct/USB, Micro-USB data cable, \texttt{python ep\_teleop\_sdk.py --conn rndis --duration 30 --out logs/rndis\_run.csv}.
\textbf{Result.} Complete 30\,s run: 299 rows over 29.9\,s (10\,Hz) with battery, attitude and all four wheel speeds logged and no missing values; the newest ESC packet arrived 0.04\,s before the file ended, so the link stayed up for the whole run. Forward and lateral commands were used ($|v_x|=|v_y|=0.3$\,m/s), yaw was not. The telemetry arrival interval was $102.3\pm8.8$\,ms (min 98.3, max 202.5, 95th percentile 107.0, $n=291$); the single 200\,ms gap is one late packet. (An earlier wired attempt was cut to 9.6\,s by a USB enumeration failure and was replaced by this run.)
\todo{attach terminal snapshot of the live telemetry display}

\subsection{(iii) Wireless AP run (30\,s)}
Switch on AP mode, laptop joined to the robot's access point (network profile \texttt{RMEP-21b442}), \texttt{python ep\_teleop\_sdk.py --conn ap --duration 30 --out logs/ap\_run.csv}.
\textbf{Result.} 299 rows over 29.9\,s (10\,Hz); battery, attitude and wheel speeds were logged and all three motions were used (forward and lateral at 0.3\,m/s, yaw at $45^\circ$/s). The ESC telemetry arrival interval was $120.2\pm52.3$\,ms (min 17.3, max 319.6, 95th percentile 222.7, $n=247$).
\todo{attach terminal snapshot of the live telemetry display}

\subsection{(iv) Comparative latency analysis}
Ping: \texttt{scripts/ping\_bench.sh rndis} and \texttt{ap} (100 packets each). Plots by \texttt{scripts/analyze\_telemetry.py}. The telemetry interval is $\Delta t=t_k-t_{k-1}$ over ESC packet arrival times (network-level), next to the logger loop interval (host scheduling only).

\begin{figure}[H]\centering
\figorti{../figures/fig_dt.png}{0.95}{run analyze\_telemetry.py to generate fig\_dt.png}
\caption{Telemetry sampling interval $\Delta t$ for wired RNDIS and wireless AP.}
\end{figure}
\begin{figure}[H]\centering
\figorti{../figures/fig_ping.png}{0.95}{run analyze\_telemetry.py to generate fig\_ping.png}
\caption{Ping round-trip time, $N=100$.}
\end{figure}

\begin{table}[H]\centering\small
\begin{tabular}{@{}lrr@{}}
\toprule
 & \textbf{Wired RNDIS} & \textbf{Wireless AP}\\
\midrule
Ping loss ($N=100$) & 0\,\% & 0\,\%\\
RTT mean (ms) & 0.451 & 28.3\\
RTT std (ms) & 0.252 & 36.7\\
Jitter, mean $|{\rm RTT}_i-{\rm RTT}_{i-1}|$ (ms) & 0.137 & 32.1\\
RTT min / 95th pct / max (ms) & 0.28 / 0.56 / 2.75 & 1.68 / 127.3 / 152.0\\
Telemetry $\Delta t$ mean $\pm$ std (ms) & $102.3\pm8.8$ (n=291) & $120.2\pm52.3$ (n=247)\\
Telemetry $\Delta t$ min / max (ms) & 98.3 / 202.5 & 17.3 / 319.6\\
\bottomrule
\end{tabular}
\caption*{\textbf{Measured latency statistics.} Both telemetry runs are complete 30\,s logs.}
\end{table}

\textbf{Discussion.} Wireless AP responsiveness is clearly worse than wired, even though neither link lost any ping.
\begin{itemize}
  \item \textbf{Latency:} mean RTT rises from 0.45\,ms to 28.3\,ms, about 63 times. The wired link never exceeded 2.8\,ms, while the AP link reached 152\,ms.
  \item \textbf{Jitter:} the AP RTT has a standard deviation of 36.7\,ms (larger than its own mean), against 0.25\,ms wired, and a mean successive difference of 32.1\,ms against 0.14\,ms. The distribution is heavy-tailed: the 95th percentile is 127.3\,ms against 0.56\,ms.
  \item \textbf{Telemetry intervals:} wired packets arrive almost exactly every 100\,ms ($102.3\pm8.8$\,ms, with one 202\,ms outlier). Over AP the stream is bursty (Figure~1): gaps of about 200 to 320\,ms are followed by very short intervals, down to 17\,ms, as delayed packets are delivered together. The mean stays near the nominal period (120\,ms) while the standard deviation is about six times larger. The logger saw 248 distinct ESC packets in 299 ticks, so stale values were repeated during the delays.
  \item \textbf{Command lag:} a command takes one-way delay plus up to 50\,ms for the 20\,Hz dispatch. With RTT spikes up to 150\,ms over AP, a key release or Space press can reach the robot about 0.15\,s late, which at 0.3\,m/s is roughly 4.5\,cm of extra travel.
\end{itemize}
Further points:
\begin{itemize}
  \item \textbf{Responsiveness:} wired USB is a point-to-point link with no contention, so RTT is low and tight. A 2.4/5\,GHz AP adds medium-access contention, retransmissions and power-save wake-ups, which raises the mean and, more importantly, the tail.
  \item \textbf{Command lag:} a command reaches the robot after one-way delay plus the 20\,Hz dispatch period (up to 50\,ms) and the 100\,ms telemetry period bounds how stale the feedback is. Over AP the one-way delay variance can exceed the dispatch period, so lag becomes visibly irregular.
  \item \textbf{Jitter:} appears as spread in $\Delta t$: a 10\,Hz stream should give $\Delta t\approx100$\,ms; late packets followed by bursts give a heavy right tail and a compensating short interval.
  \item \textbf{Safety link:} a missed stop packet over AP is the practical risk, which is why the Space key and release timeout matter more on the wireless link.
\end{itemize}

\section{Problem 3: Wireless teleoperation and telemetry in ROS 2}

\subsection{(i) Driver over hotspot}
\begin{Verbatim}[fontsize=\small]
source ~/rm_setup.sh
ros2 launch robomaster_ros ep.launch conn_type:=ap
\end{Verbatim}
Verification that the node initialises, advertises \texttt{/cmd\_vel}, and publishes \texttt{/odom} and \texttt{/joint\_states}:
The driver log reports \texttt{Try to connect via ap to robot with sn None} followed by \texttt{Connected}. Running nodes: \texttt{/robomaster}, \texttt{/robot\_state\_publisher}, \texttt{/joint\_state\_publisher}. \texttt{/cmd\_vel} has one subscriber (\texttt{/robomaster}), and \texttt{/odom}, \texttt{/joint\_states}, \texttt{/imu}, \texttt{/battery}, \texttt{/tf} are advertised. Measured rates (8\,s windows, \texttt{logs/ap\_ros\_verification.txt}):
\begin{itemize}
  \item \texttt{/odom}: 9.8\,Hz average, with intervals between 4\,ms and 226\,ms (std 46 to 51\,ms).
  \item \texttt{/joint\_states}: 10.2\,Hz average, intervals between 19\,ms and 101\,ms (std 12 to 14\,ms).
\end{itemize}
Both topics average the configured 10\,Hz. The \texttt{/odom} intervals spread much wider (4 to 226\,ms), consistent with odometry packets arriving over Wi-Fi in bursts as in Problem 2.

\subsection{(ii) Keyboard pipeline}
\begin{Verbatim}[fontsize=\small]
ros2 run teleop_twist_keyboard teleop_twist_keyboard \
    --ros-args -r /cmd_vel:=/cmd_vel
\end{Verbatim}
The run lasted about 289\,s with 81 \texttt{/cmd\_vel} messages (maximum $|v_x|=0.55$\,m/s, $|\omega_z|=1.0$\,rad/s), with roughly 21\,s of forward driving and 40\,s of rotation; the net heading change is $+1447^\circ$ (about four turns), from repeated in-place turns at the corners of the path. \textbf{No lateral strafing was commanded} ($\max|v_y|=0$ in the bag), so the strafing requirement of this step is not yet met. \todo{short repeat run with Shift+J / Shift+L strafing legs}

\subsection{(iii) Telemetry acquisition}
\begin{Verbatim}[fontsize=\small]
ros2 bag record -o logs/bag_ap /odom /joint_states /cmd_vel
\end{Verbatim}
Result (\texttt{logs/ap\_bag\_info.txt}): 289.5\,s, 5804 messages, MCAP storage: \texttt{/odom} 2827 (9.8\,Hz), \texttt{/joint\_states} 2896 (10.0\,Hz), \texttt{/cmd\_vel} 81. The largest gap between consecutive \texttt{/odom} messages is 0.26\,s and only 3 gaps exceed 0.25\,s, so recording was continuous with no network stall.

\subsection{(iv) State estimation and odometry trajectory}
\begin{figure}[H]\centering
\figorti{../figures/fig_odom.png}{0.6}{run plot\_odom.py on the bag to generate fig\_odom.png}
\caption{Estimated 2-D trajectory from \texttt{/odom}.}
\end{figure}
\textbf{Endpoint comparison.} Expressed in the start pose frame ($x$ forward, $y$ left), the estimated end is $(X,Y)=(0.881,\,0.447)$\,m after an integrated path length of 14.86\,m. \todo{tape-measured true end position (same frame) and the resulting error in m and as \% of path length: not measured for this run}

\textbf{Sources of drift.}
\begin{itemize}
  \item \textbf{Dead-reckoning integration:} pose is the running integral $\mathbf p_{k+1}=\mathbf p_k+\mathbf R(\psi_k)\mathbf v_k\Delta t$ of measured wheel velocities. Each step's error is never corrected, so position error grows without bound, and heading error is the worst contributor because it rotates every later displacement: a heading error $\delta\psi$ yields a lateral position error of about $d\,\delta\psi$ after travelling distance $d$.
  \item \textbf{Mecanum roller slip:} the lateral and yaw motions rely on rollers sliding at 45$^\circ$, so wheel rotation is only loosely tied to ground displacement. Strafing and rotating slip much more than driving straight, and the effective wheel radius changes with the contact point on the rollers. Wheel encoders measure the wheel, not the ground, so slip is invisible to odometry.
  \item \textbf{Model and discretisation errors:} the nominal $R$, $l_x$, $l_y$ differ from the real values (tyre wear, load), and 10\,Hz sampling with first-order integration ignores motion within a step during turns. Accelerations cause transient slip too.
  \item \textbf{Sensing noise:} gyro bias and encoder quantisation add a slowly growing error, and any late or dropped \texttt{/odom} packet over Wi-Fi coarsens the integration step.
\end{itemize}

\section{Problem 4: Theory and diagnostics}

\subsection{(i) Network topology characterisation}
Qualitative rows follow from the topologies; measured rows come from the runs above.

\begin{table}[H]\centering\small\raggedright
\begin{tabular}{@{}p{3.2cm}p{5.6cm}p{5.6cm}@{}}
\toprule
\textbf{Metric} & \textbf{Wired USB (\texttt{rndis})} & \textbf{Direct hotspot (\texttt{ap})}\\
\midrule
Host interface name & \texttt{enx52a542fd448f} & \texttt{wlp1s0}\\
Host IP and netmask & 192.168.42.3/24 & 192.168.2.32/24\\
Robot endpoint IP & 192.168.42.2 & 192.168.2.1\\
Round-trip latency (avg $\pm$ std, $N=100$) & $0.45\pm0.25$\,ms (jitter 0.14\,ms) & $28.3\pm36.7$\,ms (jitter 32.1\,ms)\\
Host internet retained? & Yes (observed: the Wi-Fi stayed connected while the USB link was up). & No (observed: the AP profile has no gateway and the only Wi-Fi radio is tied to the robot).\\
Multi-robot scalability & Poor. One cable and one USB interface per robot, and every robot uses the same 192.168.42.2 subnet, so they collide on one host without extra isolation. & Poor. A host Wi-Fi adapter joins one AP at a time and every robot AP uses 192.168.2.1, so one adapter serves one robot.\\
\bottomrule
\end{tabular}
\caption{Physical connection topology comparison.}
\end{table}

\subsection{(ii) Mecanum kinematic derivation}
\textbf{Setup.} Body frame: $x$ forward, $y$ left, $\omega_z$ counter-clockwise (as in Figure~1 of the assignment). Wheel $i$ has radius $R$ and angular velocity $\omega_i$ (positive drives the vehicle forward) and sits at $(x_i,y_i)$ with $x_i=\pm l_x$, $y_i=\pm l_y$:
\[
\text{FL}=(+l_x,+l_y),\quad \text{FR}=(+l_x,-l_y),\quad \text{RL}=(-l_x,+l_y),\quad \text{RR}=(-l_x,-l_y).
\]
Rollers sit at angle $\gamma_i=\pm45^\circ$ to the wheel's driving direction (the standard RoboMaster layout: FL and RR have one roller handedness, FR and RL the other). The figure does not draw rollers, so this is the assumed layout.

\textbf{Step 1: velocity of the contact point.} For a rigid body twist $\mathbf v_b=[v_x,v_y,\omega_z]^T$ the velocity of the point at $(x_i,y_i)$ is
\[
\mathbf v_i=\begin{bmatrix}v_x-\omega_z y_i\\ v_y+\omega_z x_i\end{bmatrix}.
\]

\textbf{Step 2: wheel and roller constraint.} The contact point moves with the rim speed $u_i=R\omega_i$ along the wheel's driving direction $\hat{\mathbf x}$, plus a free sliding speed $s_i$ along the roller axis $\mathbf r_i=(\cos\gamma_i,\sin\gamma_i)$:
\[
\mathbf v_i=u_i\hat{\mathbf x}+s_i\mathbf r_i .
\]
The $y$ component gives $s_i=v_{iy}/\sin\gamma_i$. Substituting into the $x$ component:
\[
R\omega_i=u_i=v_{ix}-v_{iy}\cot\gamma_i,\qquad \sigma_i:=\cot\gamma_i=\pm1 .
\]

\textbf{Step 3: insert the contact velocity.}
\[
R\omega_i=v_x-\sigma_i v_y-\omega_z\,(y_i+\sigma_i x_i).
\]
With $\sigma=(+1,-1,-1,+1)$ for (FL, FR, RL, RR), every term $y_i+\sigma_i x_i$ equals $\pm(l_x+l_y)$. Define $k=l_x+l_y$:
\[
\begin{aligned}
\omega_{fl}&=\tfrac1R\,(v_x-v_y-k\,\omega_z), &
\omega_{fr}&=\tfrac1R\,(v_x+v_y+k\,\omega_z),\\
\omega_{rl}&=\tfrac1R\,(v_x+v_y-k\,\omega_z), &
\omega_{rr}&=\tfrac1R\,(v_x-v_y+k\,\omega_z).
\end{aligned}
\]

\textbf{Inverse kinematics.}
\[
\mathbf J_{IK}=\frac1R\begin{bmatrix}1&-1&-k\\1&\;\;1&\;\;k\\1&\;\;1&-k\\1&-1&\;\;k\end{bmatrix}
\]

\textbf{Forward kinematics.} $\mathbf J_{IK}$ is $4\times3$ with full column rank, so the least-squares inverse is the pseudo-inverse $\mathbf J_{FK}=(\mathbf J_{IK}^T\mathbf J_{IK})^{-1}\mathbf J_{IK}^T$. Its three columns $\mathbf c_1=(1,1,1,1)$, $\mathbf c_2=(-1,1,1,-1)$, $\mathbf c_3=(-k,k,-k,k)$ (up to the factor $1/R$) are mutually orthogonal: $\mathbf c_1\!\cdot\!\mathbf c_2=0$, $\mathbf c_1\!\cdot\!\mathbf c_3=0$, $\mathbf c_2\!\cdot\!\mathbf c_3=k+k-k-k=0$. Hence
\[
\mathbf J_{IK}^T\mathbf J_{IK}=\frac1{R^2}\operatorname{diag}(4,\,4,\,4k^2),
\]
which gives
\[
\mathbf J_{FK}=\frac R4\begin{bmatrix}1&1&1&1\\-1&1&1&-1\\-\tfrac1k&\tfrac1k&-\tfrac1k&\tfrac1k\end{bmatrix},\qquad
\begin{aligned}v_x&=\tfrac R4(\omega_{fl}+\omega_{fr}+\omega_{rl}+\omega_{rr})\\
v_y&=\tfrac R4(-\omega_{fl}+\omega_{fr}+\omega_{rl}-\omega_{rr})\\
\omega_z&=\tfrac R{4k}(-\omega_{fl}+\omega_{fr}-\omega_{rl}+\omega_{rr}).\end{aligned}
\]
\textbf{Check.} $\mathbf J_{FK}\mathbf J_{IK}=\mathbf I_3$. For example, the $(1,1)$ entry is $\tfrac R4\cdot\tfrac1R(1+1+1+1)=1$ and the $(3,3)$ entry is $\tfrac R{4k}\cdot\tfrac1R(k+k+k+k)=1$. $\mathbf J_{IK}\mathbf J_{FK}\neq\mathbf I_4$: four wheels span only a 3-D twist space, and the left-over direction $(1,-1,-1,1)$ is a wheel-speed combination that produces no body motion, which is why wheels fighting each other show up as slip.

\textbf{Sanity checks.} Forward: all wheels $+v_x/R$. Strafe left ($v_y>0$): FL and RR spin backward, FR and RL forward. Rotation ($\omega_z>0$, CCW): left wheels backward, right wheels forward. All match the known mecanum behaviour.

\textbf{Numbers.} Nominal RoboMaster EP values (confirm against the datasheet before quoting): $R=0.05$\,m, $l_x=l_y=0.10$\,m, so $k=0.20$\,m.
Conversion to the SDK convention ($y$ right, $\omega_z$ clockwise): $v_y^{\rm SDK}=-v_y$, $\omega_z^{\rm SDK}=-\omega_z$.

\subsection{(iii) ROS 2 coordinate transform inspection}
\texttt{ros2 run tf2\_tools view\_frames} was run with the driver up.
\begin{figure}[H]\centering
\figorti{../figures/frames.pdf}{1.0}{attach frames.pdf generated by view\_frames and save it as figures/frames.pdf}
\caption{TF tree generated by \texttt{tf2\_tools view\_frames} (vector PDF, zoom in for detail).}
\end{figure}
Parent and child links, read from the generated tree (\texttt{figures/frames.pdf}) and \texttt{ros2 topic info /tf -v}, which lists exactly two publishers of \texttt{/tf}:
\begin{itemize}
  \item \texttt{odom} $\to$ \texttt{base\_link}: the only dynamic link at the top of the tree, published at about 10\,Hz by the driver node \texttt{/robomaster} (its chassis module, from the dead-reckoned pose). \texttt{odom} is the root of the tree.
  \item \texttt{base\_link} $\to$ \texttt{chassis\_base\_link}: a fixed (static) link from the robot description, published by \texttt{/robot\_state\_publisher}.
  \item \texttt{chassis\_base\_link} $\to$ \texttt{arm\_base\_link}: a fixed (static) link (\texttt{arm\_base\_joint}, offset $(0,\,0.0010,\,0.0906)$\,m), also published by \texttt{/robot\_state\_publisher}.
\end{itemize}
Beyond these, \texttt{/robot\_state\_publisher} also publishes the moving links, driven by \texttt{/joint\_states}: the four wheel links (children of \texttt{chassis\_base\_link}) and the arm and gripper chain (\texttt{arm\_base\_link} $\to$ \texttt{arm\_1\_link} $\to$ \texttt{arm\_2\_link} $\to$ \ldots), each at about 10\,Hz. Sensor, LED and camera frames hang off \texttt{chassis\_base\_link} and \texttt{arm\_2\_link} as fixed links.

\section*{Deliverables}
\begin{itemize}
  \item D1: this report. D2: \texttt{ep\_teleop\_sdk.py}, \texttt{rm\_setup.sh}, \texttt{scripts/}, raw CSV logs, ROS bag. D3: video link \todo{unlisted YouTube/cloud link, $\le$45\,s}.
\end{itemize}

\end{document}
